%% =====================================================================
%%  T-23-05  The Five-Second Window
%%  -------------------------------------------------------------------
%%  One idea per file. Core is mandatory. Slots in canonical order.
%%  Max 700 words. No layout commands. No filler. New source material
%%  for this entry goes in sources/B7/T-23-05.tex -- never by
%%  rewriting this file (docs/RULES.md R0.1).
%%  =====================================================================
%% @kind: topic
%% @id: T-23-05
%% @title: The Five-Second Window
%% @subtitle: Watch the answer, not the person --- behaviour is evidence only right after a question
%% @chapter: c23
%% @order: 50
%% @status: stable
%% @sources: B7, B8
%% @updated: 2026-09-20

\topic{T-23-05}{The Five-Second Window}{Watch the answer, not the person --- behaviour is evidence only right after a question}

\begin{core}
B7's discipline, built from interrogation practice: there is no human lie detector, and reading someone's general manner is guessing. The instrument is question-locked observation --- deliver the question (the stimulus), then look and listen at once for the first deceptive sign; anything inside the following five seconds can be attributed to the question, and only a cluster of two or more indicators counts. One behaviour is noise; timing is what turns manner into evidence.
\end{core}
\begin{mechanism}
Global assessment --- is he shifty, is she calm? --- fails because it must guess why a behaviour occurs; habits, nerves and baseline manner all masquerade as signals. Timing fixes the attribution: speech runs at roughly 125--150 words a minute while thought runs an order of magnitude faster, so a mind answering honestly has already moved on to content, while a mind managing a deception is still working the question --- hence the five-second span, beyond which a behaviour increasingly belongs to something other than the question. The second rule is the cluster: any two or more deceptive indicators, in any mix of word and body, form the reading, because individuals have idiosyncratic habits (one gesture means nothing) and because confidence scales with the count. Both channels must be captured at once --- B7's L-squared discipline, dual attention the brain resists and must be trained into --- because the catalogue spans words (nonanswers, qualifiers, selective memory) and body (a behavioural pause, hiding mouth or eyes, a pre-answer throat-clear, hand-to-face activity, grooming gestures). And truthful responses carry their own signature: a reader who watches only the deceptive list will eventually read nerves as guilt. The measured ceiling over the whole method is filed under T-23-09.
\end{mechanism}
\begin{conditions}
Requires a live exchange you can steer --- a question you choose and a subject you can observe unhurried. It degrades on broadcast or written channels, against subjects whose baseline you cannot sample, and under time pressure that collapses the window. It never reaches certainty: B7 is explicit that clusters raise confidence, not proof, and that single-sign reads (the microexpression school of television) produce false positives at scale.
\end{conditions}
\begin{application}
  \item Write the question before the meeting and deliver it clean. A muddled stimulus makes every subsequent behaviour unattributable.
  \item After the question lands, go silent and open both channels. Five seconds, counted --- most readers quit at two.
  \item Score the cluster, not the gesture: two indicators, tentative; three-plus, act. One indicator, file it away as nothing.
  \item Re-fire the same stimulus much later in the conversation. A cluster that re-arms on the question but not on small talk is anchored to the issue.
  \item Sample the baseline first: how do they behave on questions you know are innocent? The window reads deviations, so you need the ordinary.
\end{application}
\begin{feedback}
  \item Working: your reads arrive with a count --- \emph{three indicators on the money question, none on the rest}. That is the method running.
  \item Working: you stop rating manner entirely. Charm and calm stop feeling like evidence.
  \item Failing: every nervous person reads as lying. You are scoring singles and skipping the baseline; the cluster rule exists precisely for you.
  \item Failing: you confront on a cluster of two with no re-test. The window supports suspicion; only behaviour over time supports confrontation.
\end{feedback}
\begin{failure}
The method's failure mode is the false positive with a certificate: a person trained to five seconds but not to baselines converts everyone's quirks into clusters, and now has a vocabulary for prejudice. B7's own guard is humility --- no human lie detector, confidence proportional to cluster size, truthful signatures matter as much as deceptive ones --- and the moment the guard is dropped, the instrument becomes an accusation machine.
\end{failure}
\begin{countermeasures}
  \item When someone runs the window on you, keep answers boring: answer the question asked, once, in plain words. The catalogue feeds on elaboration.
  \item Notice interviewers who go quiet right after a question --- the silence is the instrument's intake stroke; do not fill it nervously.
  \item For your own people-reading: demand a stimulus and a count before you believe any read, yours included.
  \item The social obstacle --- nobody would lie to me --- is the window's true enemy; it keeps the question from ever being asked (T-23-04).
\end{countermeasures}

\sources{B7, B8}
\seesources
\seealso{T-23-04, T-23-06, T-06-02, T-23-09}
